\documentclass{article}
\usepackage{amsmath} 
\usepackage{amsfonts}
\usepackage{amssymb}
\usepackage[parfill]{parskip}
\usepackage{esvect}
\usepackage[thinc]{esdiff}
\usepackage{mathtools}
\usepackage{physics}
\usepackage{amsthm}
\usepackage{xfrac}
\usepackage{enumitem}

\DeclareMathOperator{\lub}{lub}
\DeclareMathOperator{\gub}{gub}

\begin{document}

\title{MATH 411: Week 8}
\author{Jacob Lockard}
\date{24 March 2026}

\maketitle

\section*{Lemma 1}

\begin{quote}
    \itshape
    Suppose $f$ is defined on an interval $(a, b)$ with $f(x) < M$
    for all $x \in (a,b)$, and let $p \in (a,b)$. Then,
    \begin{align*}
        L = \limsup_{x \to p^+} f(x) \leq M \,.
    \end{align*}
    A similar result holds for minima and limits inferior.
\end{quote}
Since $f$ is bounded on $(a, b)$, we already know $L$
is finite. Suppose $L > M$ and let $\epsilon = L-M > 0$.
Then if there's an $x \in (p,b) \subseteq (a,b)$ such that
$f(x) \in (L - \epsilon, L + \epsilon)$, we have
\begin{align*}
    f(x) &> L - \epsilon = L - (L-M) = M \,.
\end{align*}
Since we know $f(x) < M$, no such $x$ exists.
We've shown that there's a neighborhood around $L$ 
containing no outputs of $f$ corresponding to inputs
to the right of $p$. But $L$ is a limit superior,
so we conclude $L \leq M$.\qedsymbol

\section*{Setup}

For Problems 8.21--24, assume that
\begin{quote}
    \itshape
    Let $f$ and $g$ be real, differentiable functions on some open interval $I$
    with $a \in I$, and $g'(x) \neq 0$ for all $x \in I$.
    Suppose
    \begin{align*}
        \lim_{x \to a^+} \frac{f'(x)}{g'(x)} = A \,.
    \end{align*}
    Also suppose $\lim_{x \to a^+} f(x) = 0$ and $\lim_{x \to a^+} g(x) = 0$.
\end{quote}

\section*{Problem 8.21}

\begin{quote}
    \itshape
    Let $p$ and $q$ be any real numbers such that $p < A < q$,
    and choose an $r$ and $s$ such that $A < r < q$ and $p < s <A$.
    Let $b \in I$ with $a < b$. Explain why there exists some $c \in (a,b)$
    such that for all $x$ with $a < x < c$ we have
    \begin{align*}
        s < \frac{f'(x)}{g'(x)} < r \,.
    \end{align*}
\end{quote}

Let $\epsilon = \min\{ r-A, A-s \}$, which is positive since
$r > A$ and $A > s$.
Then there's a $\delta > 0$
such that if $0 < |x-a| < \delta$, then
\begin{gather*}
    \Bigl| \frac{f'(x)}{g'(x)} - A \Bigr| < \epsilon \\
    -\epsilon < \frac{f'(x)}{g'(x)} - A < \epsilon \\
    A-\epsilon < \frac{f'(x)}{g'(x)} < A+\epsilon \,,
\end{gather*}
which we can rewrite:
\begin{align*}
    s = A - (A - s) \leq A-\epsilon < \frac{f'(x)}{g'(x)} < A+\epsilon \leq A + (r-A) = r \,.
\end{align*}
Let $\gamma = \min\{ b-a, \delta \}$ and let $c \in (a, a + \gamma)$.
$\gamma$ is positive, since $b > a$ and $\delta > 0$, so $c$ exists.
Further, $c \in (a, b)$, since
\begin{align*}
    c < a + \gamma \leq a + (b - a) = b \,.
\end{align*}
If $x \in (a, c)$, then
\begin{gather*}
    a < x < c < a + \gamma \leq a + \delta \\
    0 < x - a < \delta \\
    0 < |x-a| < \delta \,,
\end{gather*}
and thus
\begin{align*}
    s < \frac{f'(x)}{g'(x)} < r \,.
\end{align*}
So we've found a $c$ with the desired properties.
\qedsymbol

\section*{Problem 8.22}

\begin{quote}
    \itshape
    Let $x, y \in (a,c)$ with $x < y$. Explain why there exists $t \in (x, y)$
    such that
    \begin{align*}
        s < \frac{f(y) - f(x)}{g(y) - g(x)} = \frac{f'(t)}{g'(t)} < r \,.
    \end{align*}
\end{quote}

Since $I$ is an open interval and both $x$ and $y$ are in $I$, neither
$x$ nor $y$ are the endpoints of $I$, which ensures $f$ and $g$
are defined and differentiable (and thus continuous) on $[x, y]$.
So the Cauchy mean value theorem ensures there's a $t \in (x, y)$
such that
\begin{align*}
    \bigl( f(y) - f(x) \bigr) g'(t) &= \bigl( g(y) - g(x) \bigr) f'(t) \,.
\end{align*}
If $g(y) - g(x) = 0$, then the mean value theorem implies that
there's a $p \in (x, y)$ such that
\begin{align*}
    (y-x)g'(p) &= 0 \\
    g'(p) &= 0 \,,
\end{align*}
where the last equality is justified since $x \neq y$.
However, by assumption $g'(p) \neq 0$, so we conclude
$g(y)-g(x) \neq 0$. Also, $g'(t) \neq 0$. Now we can write
\begin{align*}
    \frac{f(y) - f(x)}{g(y) - g(x)} &= \frac{f'(t)}{g'(t)} \,,
\end{align*}
as desired. \qedsymbol

\section*{Problem 8.23}

\begin{quote}
    \itshape
    Prove that for all $y \in (a,c)$,
    \begin{align*}
        p < s \leq \frac{f(y)}{g(y)} \leq r < q \,.
    \end{align*}
\end{quote}

Let $y \in (a, c)$, and consider
\begin{align*}
    \lim_{x \to a^+} \frac{f(y) - f(x)}{g(y) - g(x)}
    &=  \frac{\lim_{x \to a^+} \bigl( f(y) - f(x) \bigr)}{\lim_{x \to a^+} \bigl(g(y) - g(x)\bigr)} \\
    &=  \frac{f(y) - \lim_{x \to a^+} f(x) }{ g(y) - \lim_{x \to a^+} g(x)} \\
    &= \frac{f(y)}{g(y)} \,,
\end{align*}
where the equalities are justified by the fact that
$f$ and $g$ are continuous at $x$.
Let $\alpha \in I$ with $\alpha < a$ and let $\beta \in (y, c)$.
$\alpha$ exists, since $I$ is an open interval and $a \in I$.
Since $(\alpha, \beta) \subseteq I$, both $f$ and $g$ are defined on
$(\alpha, \beta)$.
The previous problem ensures that for any $x \in (\alpha, \beta)$, we have
\begin{align*}
    s < \frac{f(y) - f(x)}{g(y) - g(x)} < r \,.
\end{align*}
Lemma 1 now ensures
\begin{align*}
    s \leq \liminf_{x \to a^+} \frac{f(y) - f(x)}{g(y) - g(x)}
    = \lim_{x \to a^+} \frac{f(y) - f(x)}{g(y) - g(x)}
    = \limsup_{x \to a^+} \frac{f(y) - f(x)}{g(y) - g(x)}
    \leq r \,.
\end{align*}
\qedsymbol

\section*{Problem 8.24}

\begin{quote}
    \itshape
    Conclude the proof of the theorem, that is that
    \begin{align*}
        \frac{f(x)}{g(x)} \to A \text{ as } x \to a \text{ (from the right). }
    \end{align*}
\end{quote}

The previous problem ensures that
\begin{align*}
    p < \frac{f(y)}{g(y)} < q \,.
\end{align*}
So Lemma 1 ensures
\begin{align*}
    p \leq \liminf_{x \to a^+} \frac{f(y)}{g(y)} \leq \limsup_{x \to a^+} \frac{f(y)}{g(y)} \leq q \,.
\end{align*}
Since $p$ and $q$ are arbitrary real numbers with $p < A < q$, we conclude that
\begin{align*}
    \liminf_{x \to a^+} \frac{f(y)}{g(y)} = A =  \limsup_{x \to a^+} \frac{f(y)}{g(y)} \,.
\end{align*}
\qedsymbol

\section*{Problem 10.2}

\begin{quote}
    \itshape
    Consider the sequence of functions $\{ g_n \}$ on the set
    $[0,1]$ where $g_n(x) = x^n$. (Note that each $g_n$ is continuous
    on $[0,1]$.)

    \begin{enumerate}[label=(\alph*)]
        \item Show that $\{g_n\}$ converges pointwise and find the limit
        function $g$.
        \item Is $g$ continuous on $[0,1]$?
    \end{enumerate}
\end{quote}

\textbf{(a)} Let $x \in [0,1]$.
If $x = 1$, then the sequence $\{x^n\}$ clearly converges to $1$.
On the other hand, if $x < 1$, then the sequence
$\{x^n\}$ converges to $0$. So,
\begin{align*}
    g = \lim_{n\to\infty} g_n(x) = \begin{cases}
        1 \,, \quad x = 1 \,, \\
        0 \,, \quad x < 1 \,.
    \end{cases}
\end{align*}
We've shown that $\{g_n\}$ converges pointwise to $g$.

\textbf{(b)} Clearly not.

\end{document}